\documentclass[12pt]{article}
\usepackage[T1]{fontenc}
\usepackage[utf8]{inputenc}
\usepackage{lmodern}
\usepackage{textcomp}
\usepackage{enumitem}
\usepackage{geometry}
\geometry{margin=1in}
\begin{document}
\begin{center}
Answer Key - Sociology - Graduate Level Assessment 23 \
\small (Answers are ordered exactly as in the question paper.)
\end{center}

\section*{Multiple Choice}
\begin{enumerate}[label=\arabic*.]
\item \textbf{Answer:} C
\\ \textit{Options:} A) they have a deep-rooted sense of ineffectiveness and are striving for control; B) they are in pursuit of the 'perfect' body, fuelled by images of beauty; C) they perceive a part of their body as stigmatizing, in relation to a cultural ideal; D) their male partners pressurize them to look like supermodels
\\ \textit{Note:} The correct answer highlights the sociocultural pressures women face, where they may view certain body parts as stigmatizing when compared to prevailing cultural ideals of beauty, prompting them to seek cosmetic surgery as a means of aligning with those ideals.
\item \textbf{Answer:} C
\\ \textit{Options:} A) the mass production of standardized products for passive audiences; B) television based on producer-broadcaster rather than publisher-broadcaster models; C) a diverse range of products aimed at niche markets; D) increasing numbers of advertisements for motoring and car-related products
\\ \textit{Note:} Post-Fordist forms of media production and consumption are characterized by a shift from mass production to a focus on customization and diversity, targeting specific niche markets to cater to varied consumer preferences and cultural identities.
\item \textbf{Answer:} D
\\ \textit{Options:} A) the privatization of public services by the Conservative government; B) the lifestyle practice of shopping in peer groups; C) the form of tuberculosis suffered by those who collect stamps; D) the provision of health, housing, and education services by the state
\\ \textit{Note:} The term 'collective consumption' refers to the state's role in providing essential services such as health, housing, and education, highlighting the social and economic structures that facilitate communal access to these resources.
\item \textbf{Answer:} C
\\ \textit{Options:} A) British-born second generation immigrants from the Asian subcontinent; B) White Americans who wanted to convert to Islam; C) African-Americans who felt excluded from the 'ethnic melting pot' in the USA; D) African-Caribbeans who lived in the inner cities and had a distinctive youth culture
\\ \textit{Note:} The Nation of Islam appealed to African-Americans who felt excluded from the 'ethnic melting pot' in the USA by promoting a distinct identity and cultural pride that countered systemic racism and offered a sense of empowerment and community.
\item \textbf{Answer:} A
\\ \textit{Options:} A) active non-work and independence after retirement; B) full time employment, family-building and adult responsibility; C) illness, isolation and increasing dependence on others; D) the transition from education to work, and distinctive youth cultures
\\ \textit{Note:} The 'third age' of the life course is characterized by a phase of active engagement in leisure and personal pursuits, emphasizing independence and self-fulfillment after the cessation of formal employment.
\end{enumerate}

\section*{True / False}
\begin{enumerate}[label=\arabic*.]
\setcounter{enumi}{5}
\item \textbf{Answer:} True
\\ \textit{Options:} A) True; B) False
\\ \textit{Note:} Bourdieu's concept of cultural capital highlights how individuals inherit and accumulate knowledge, skills, and cultural assets that influence their social mobility and reinforce social stratification, thereby facilitating the reproduction of social class across generations.
\item \textbf{Answer:} True
\\ \textit{Options:} A) True; B) False
\\ \textit{Note:} According to Marx, the development of the labor movement is essential for the working class to gain class consciousness and unite in pursuit of their collective interests, thereby transforming them from a "class in itself" (a group with shared economic conditions) to a "class for itself" (a politically organized group advocating for their rights and interests).
\item \textbf{Answer:} False
\\ \textit{Options:} A) True; B) False
\\ \textit{Note:} During the nineteenth century, homosexuality was largely stigmatized and criminalized, leading to widespread societal discrimination and repression rather than celebration or pride among gay individuals.
\item \textbf{Answer:} True
\\ \textit{Options:} A) True; B) False
\\ \textit{Note:} The 1990 reform of the National Health Service introduced the concept of self-governing NHS trusts, which allowed hospitals to operate more autonomously and engage in competition for contracts with health authorities, fundamentally shifting the structure of healthcare delivery in the UK.
\item \textbf{Answer:} False
\\ \textit{Options:} A) True; B) False
\\ \textit{Note:} A document can be considered authentic based on its origin and integrity, but it may still contain inherent biases or subjective interpretations, thus making the criteria of being free from political bias insufficient for determining authenticity.
\end{enumerate}

\section*{Long Form Response}
\begin{enumerate}[label=\arabic*.]
\setcounter{enumi}{10}
\item \textbf{Answer:} Metropolitan Statistical Areas (MSAs) are defined primarily by their inclusion of at least one central city with a population of 50,000 or more, alongside surrounding densely populated urbanized counties. This delineation is crucial as it encapsulates the interconnectedness of urban, suburban, and rural areas, allowing for a comprehensive analysis of development patterns. The significance of MSAs lies in their utility for sociological research, as they facilitate the examination of demographic trends, economic activity, and social dynamics across diverse geographic landscapes. By understanding the characteristics and boundaries of MSAs, researchers can better analyze how urbanization influences rural and suburban transitions, ultimately shaping policies and planning efforts that address regional growth and development disparities.
\\ \textit{Note:} The answer correctly identifies that MSAs are defined by central cities with significant populations and nearby urbanized areas, which enables a nuanced understanding of how urban, suburban, and rural development patterns are interconnected and informs sociological research.
\item \textbf{Answer:} Globalization has significantly altered the power dynamics of nation-states, often challenging their authority through various mechanisms. The proliferation of transnational corporations, international organizations, and digital communication has led to a diffusion of power away from traditional state actors. As economic and cultural exchanges transcend borders, nation-states find their regulatory capabilities undermined, particularly in areas like trade, immigration, and environmental policy. Furthermore, the rise of global civil society and non-governmental organizations empowers individuals and groups to influence policy outcomes, thereby diminishing the state's monopoly on authority. Ultimately, while globalization can extend the reach of nation-states by integrating them into a broader global framework, it simultaneously presents challenges that may weaken their sovereignty and control over domestic affairs.
\\ \textit{Note:} The answer correctly identifies that globalization disperses power from nation-states to transnational entities and influences, thereby challenging their regulatory authority and control over key policy areas.
\end{enumerate}

\vfill
\noindent\textit{End of Answer Key}
\end{document}